% Intermediate critical-distinctions matrix (outline synthesis). May later be condensed or redrawn as a figure panel.
\begin{table}[tbp]
\centering
\footnotesize
\singlespacing
\caption{Critical distinctions preserved throughout the review. Confusing either side of a row collapses evidence that the framework keeps separate.}
\label{tab:summary-distinctions}
\begin{tabular}{@{}p{2.0in}p{4.2in}@{}}
\toprule
\textbf{Distinction} & \textbf{Why it matters} \\
\midrule
Natural vs artificial sound & Tests significance but confounds acoustic differences \\
\addlinespace[2pt]
Conspecific alarm vs predator vocalization & Similar outcomes need not share recognition mechanisms \\
\addlinespace[2pt]
Na\"ive vs field-experienced receiver & Threat recognition may require experience \\
\addlinespace[2pt]
Laboratory vs field context & Foraging and exposure costs are often absent in the lab \\
\addlinespace[2pt]
Startle vs defensive avoidance & Short-latency reflexes need not predict sustained deterrence \\
\addlinespace[2pt]
Dopamine vs behavioral meaning & DA depends on region, timing, context, and valence \\
\addlinespace[2pt]
Detection vs interpretation vs action & Neural/physiological separation can occur without behavior \\
\addlinespace[2pt]
Caller identity vs reliability & Identity is widespread; reliability learning is sciurid-only so far \\
\addlinespace[2pt]
Habituation vs learned unreliability & Same declining curve; different mechanisms and interventions \\
\addlinespace[2pt]
Unpredictable timing vs longer recovery & Must be dissociated experimentally \\
\bottomrule
\end{tabular}
\end{table}
